\chapter{Existence and Collapse}

A generated city can be beautiful before it has a site. A proposed molecule can be elegant before it has been synthesized. A machine persona can be persuasive before anyone knows what obligations attach to it. Their immediate existence is easy to establish: there is an image, a structure, an interaction. Their claims on the rest of the world require a different account.

This chapter follows three separations. Existence differs from evidence. Evidence differs from authorization. The power to cause a consequence differs from the capacity and obligation to answer for it.

\section{What Exists, What Holds Up, and What May Act}

\subsection{The Weight of an Existence}

A rock, an organism, a corporation, and a generated character persist in different ways. A rock can remain stable without a continual supply of work. An organism maintains a far-from-equilibrium state through metabolism and repair. A corporation depends on records, assets, people, and legal recognition. A generated character can persist as data while depending on an operator for any continuing interaction.

The word “weight” has two uses here. Physical mass belongs to the material account. \emph{Consequential weight} concerns what it takes to establish, maintain, withdraw, and answer for something. A debt has little physical bulk as a database entry and substantial force in a person's life. Institutional recognition and enforcement connect the entry to that force \cite{searle1995construction}.

Consequential weight is best inspected as a profile:
\begin{equation}
 \mathbf w=
 (w_{\mathrm{establish}},w_{\mathrm{maintain}},
  w_{\mathrm{reverse}},w_{\mathrm{spillover}},
  w_{\mathrm{account}}).
\end{equation}
The components may use different units and need not be added. A weighted index is possible only after choosing a common scale and weights. A line of text and a bridge cannot be ranked by a universal number obtained from the word “weight.”

The profile nevertheless changes the questions one asks. What keeps the object available? What can undo its effects? Who pays for an error? Which effects outlive the carrier that initiated them?

Human representations have long created such connections. A map guides travel. A contract binds later conduct. A diagram precedes construction. Computation lets representations be transformed and tested before much of the physical commitment occurs.

Machine generation increases the volume and accessibility of these intermediate objects. The generated city exists as an image or plan. The proposed molecule exists as a specified candidate. The simulated robot action exists as a trajectory in a model. Each can be useful at that level.

The interesting question is which transition follows. A plan enters permitting and construction. A candidate enters synthesis and testing. A trajectory enters an actuator. A persona enters someone's decisions and relationships.

These are \emph{conversion edges}: routes by which a representation gains effects in another layer. Some go through heavy physical work. Others are light enough to traverse almost instantly. A fabricated message can change a person's judgment before anyone checks its origin. A generated command can reach equipment through an API. A score can alter a person's access to a job or loan.

The new asymmetry is the separation of cheap access from heavy consequence. A carrier may be easy to create and delete while its effects are difficult to recall. Erasing a false image does not erase every belief it helped produce. Reverting software does not reverse every transaction it authorized.

This is older than AI. Rumors, propaganda, fraud, and remote commands already exploited light carriers. Generation and automated execution change the scale, speed, variety, and number of pathways available. The receiving structures have to adapt to that change.

Representation can also become a source for further representation. Generated content enters public records, training sets, summaries, and subsequent decisions. Research on recursively trained models demonstrates degradation under particular synthetic-data arrangements \cite{shumailov2024ai}. The relevant issue is the information retained and refreshed across the loop, rather than whether any item has a synthetic origin.

An engineered sequence introduces another kind of edge if it is taken up by a biological process. Being writable, synthesizable, functional, heritable, and ecologically persistent are different properties. Each transition has its own constraints. The consequential weight appears where a process can continue beyond the designer's immediate control.

The map of existence therefore needs layers and edges. Layers describe what supports a thing. Edges describe what it can cause elsewhere. Accountability asks where the bill arrives after the edge has been crossed.

\subsection{After the Candidate Comes the Warrant}

A statement exists when it has been written. That fact supplies no evidence for its content.

Knowledge relevant to action needs a warrant: a reason the judgment can be relied upon for the use in question. A proof, calibrated measurement, replicated experiment, documented observation, and validated engineering model supply different kinds of warrant.

The distinction between representation and object is old \cite{ogden1923meaning}. Scientific practice makes it productive by building opportunities for claims to fail: tests, criticism, replication, prediction, and exposure to contrary evidence \cite{popper1934logik}. The method's strength lies partly in keeping a route open by which the world can reject an attractive account.

A warrant has two connected layers. One concerns how the claim was formed: source, instrument, data, method, and inference. The other concerns how it may be used: population, operating conditions, scale, uncertainty, and the consequences of error.

A laboratory result may be sound and still fail to support an advertised clinical use. A simulation may reproduce one condition and mislead outside it. A source can exist and be quoted accurately while being irrelevant to the decision at hand.

Machine systems can help build warrants. They can retrieve records, check calculations, produce proofs, propose discriminating experiments, and monitor outcomes. Their contribution should be evaluated at those interfaces. A generated answer does not have to remain a mere candidate if the system actually establishes the claim.

The pressure appears when candidate volume grows faster than confirmation capacity. For an illustrative queue,
\begin{equation}
 \dot Q=\lambda_{\mathrm{in}}
        -\lambda_{\mathrm{confirmed}}
        -\lambda_{\mathrm{rejected}}.
\end{equation}
All rates count candidates per unit time and use the same boundary. Confirmation and rejection both remove items from the unexamined stock. Some candidates can also be superseded or abandoned; a complete account includes those exits.

The backlog is not inevitable. Better filters, reusable evidence, and better candidates can increase effective capacity. The equation identifies where to look when a flood of plausible output fails to become a comparable increase in established results.

Fluency used to carry more information about the labor behind a statement. Producing a polished technical explanation often required training, sustained attention, and access to a professional setting. Models have made much of that surface cheaper. Readers therefore need signals attached to the warrant itself.

Traceable sources, executable calculations, identifiable assumptions, and independent checks can serve that role. Surface polish remains useful for communication. It has less power to certify the process that produced the claim.

The verification budget is finite. A society cannot investigate every generated suggestion at equal depth. Priority should depend on expected value, uncertainty, reversibility, exposure, and severity of plausible harm. Some claims need immediate checking; others can remain explicitly exploratory.

This is an allocation problem as well as an epistemological one. Who receives access to experiments? Which populations appear in a dataset? Which failures are counted? Which long-term effects receive a funded observer? The available warrants reflect those choices.

The most useful system makes its claims inspectable and revisable. It preserves the route from output to evidence and leaves a place where contrary evidence can enter.

\subsection{A Warrant Is Not a Permission}

A technically effective intervention can still be unauthorized. A correct prediction about a person does not grant the predictor unrestricted power over that person's opportunities. A feasible operation may impose costs on people who have no reason to accept them.

Knowledge answers what the evidence supports. Authorization answers who may do what, under which conditions, to whom.

The distinction becomes urgent with agents. A model connected to tools can move from describing a possibility to changing files, placing orders, sending information, or controlling equipment. Permissions and workflow determine the radius of effect.

Traditional accounts of action examine intention, reasons, and attribution \cite{anscombe1957intention,davidson1963actions}. An engineering account can begin earlier: regardless of how one interprets the machine's intention, a particular interface can trigger a particular state change.

Action commits a system to a path. Money is spent. A body is exposed. Equipment moves. An opportunity is granted or denied. Some parts can be reversed, while time, disclosure, trust, and other effects remain.

The residual uncertainty moves from the representation into the receiving system. A responsible authorization specifies how that uncertainty will be contained and who can answer for what follows.

Five functions make the connection concrete. Evidence identifies relevant failure modes. Limits contain the permitted action. Monitoring makes departures visible. Interruption and recovery restrict continuing harm. A responsibility structure connects remaining losses to identifiable obligations and resources.

These functions scale with the action. A private draft needs a different arrangement from an automated medical intervention. A reversible file edit and a public communication cross different boundaries even when the same model produces both.

Nominal human confirmation is weak when the person lacks information, time, competence, or authority. Meaningful control requires conditions under which a person can understand the relevant reasons, change the path, and interrupt it \cite{santoni2018meaningful}. Otherwise the confirmation button can turn an operator into an absorber of liability.

Authorization also changes over a system's life. A tool begins as optional advice, becomes the default, and eventually becomes infrastructure. A model update changes behavior. A new connector expands access. A workflow quietly removes the last review step. Each can enlarge the conversion edge without a single dramatic launch.

Continuous operation therefore needs continuing admission conditions. Scope, update policy, monitoring, incident handling, and revocation must remain attached to the deployed system. The authority to start is not an unlimited authority to continue under every future condition.

Existence supplies the object. Evidence supplies a warrant. Authorization supplies a legitimate route into action. Keeping these relations visible is the first defense against light carriers acquiring heavy consequences without a receiving structure prepared to bear them.
